% ==============================================================================
% T0 Theory / FFGFT: Document 283 (Chapter Content, DE)
% Titel: FFGFT <-> HLV: die Bruecke und die Gabelung
% Autor: Johann Pascher
% Datum: 17. Juni 2026
% Zweck: Strenge Pruefung der FFGFT<->HLV-Verbindung nach dem Drei-Teile-Kriterium.
%        Befund: gemeinsamer Z3-Zirkulant entlang Kruegers C3-Achse (Krit. 1+2),
%        fundamentale Gabelung an der Symmetrie-Ordnung (5-zaehlig/phi vs 3-zaehlig/sqrt2).
% Referenzen: Dok. 206, 230, 271, 276, 281, 282.
% Skript: ffgft_hlv_c3_bridge_probe.py
% ==============================================================================

\section*{Abstract}
Marcel Krügers HLV baut den Raum als ikosaedrischen 6D$\to$3D-Cut-and-Project-Quasikristall.
Ob FFGFT und HLV verbunden sind, wird hier streng nach dem Drei-Teile-Kriterium
geprüft: (1)~dasselbe strukturelle Objekt aus der \emph{eigenen} Geometrie jedes
Modells, (2)~eine explizite Abbildung zwischen den Herleitungen, (3)~ein
Divergenzpunkt, der bloßes Stufe-0-Gerüst ausschließt. Befund: Ein
Z$_3$-Zirkulant -- der Typ von FFGFTs Massenoperator -- lebt nachweislich in HLVs
eigener Geometrie, entlang einer C$_3$-Achse der Ikosaedergruppe (Kriterien 1+2
erfüllt). Die innere Invariante aber divergiert \emph{fundamental}: HLVs
geschützte Struktur liefert rigide $r=2/\sqrt5$ (5-zählig, $\varphi$), FFGFT
verlangt $r=\sqrt2$ (3-zählig, Koide $Q=2/3$). Kein Cut-and-Project-Parameter
erreicht $\sqrt2$ ($r(\tau)\le1$ für alle $\tau$); $\sqrt2$ taucht in HLV nur als
ungeschützter Zufallswinkel auf. Ergebnis: eine \emph{genuine} Verbindung mit
einer \emph{fundamentalen} Gabelung an der Symmetrie-Ordnung -- keine
Stufe-0-Trivialität, keine Identität. Die Gabelung sitzt im \emph{Raum}: HLVs
Phasen-Schicht ist sein Spiral-Time-Operator (nicht $\theta=0$), und auf der
Zeit/Gedächtnis-Ebene sind HLVs $\chi$-Kern und FFGFTs Dok.-282-Kern
nicht-Markovsche Kerne derselben Klasse -- ein Brücken-Kandidat (§\,8). Alle
geometrischen Aussagen sind mit
\texttt{ffgft\_hlv\_c3\_bridge\_probe.py} reproduzierbar.

% ==============================================================================
\section{Anlass und Kriterium}
% ==============================================================================

Der Austausch mit Marcel Krüger (HLV/HICE) wirft wiederholt die Frage auf, ob
FFGFT und HLV verbunden sind. Eine ehrliche Antwort braucht ein scharfes Kriterium,
sonst landet man bei Scheinverbindungen, die nur gemeinsames Gerüst teilen
(beide benutzen Planck-Einheiten, beide haben irgendwo eine Wurzel). Das
Kriterium hat drei Teile:
\begin{enumerate}
\item \textbf{Gleiches Objekt:} dasselbe strukturelle Objekt, hergeleitet aus der
      \emph{eigenen} Geometrie jedes Modells -- nicht aus geteiltem Skalengerüst.
\item \textbf{Abbildung:} eine explizite Abbildung zwischen den beiden Herleitungen.
\item \textbf{Divergenzpunkt:} ein Punkt, an dem sich die Modelle trennen --
      sonst ist die „Verbindung`` entweder triviale Identität oder Stufe-0-Gerüst.
\end{enumerate}
Dieses Dokument wendet das Kriterium konkret an. Es behauptet keine theoretische
Äquivalenz; es lokalisiert präzise, \emph{wo} sich FFGFT und HLV treffen und
\emph{wo} sie auseinanderlaufen.

% ==============================================================================
\section{Wie Marcel seinen Raum baut}
% ==============================================================================

Aus dem HLV-Homologie-Skript (Dok.~271/276) ist die Konstruktion direkt lesbar:
ein \textbf{ikosaedrischer 6D$\to$3D-Cut-and-Project-Quasikristall}. Ein
$\mathbb{Z}^6$-Gitter wird über zwei Projektoren in einen 3D-Parallelraum
(physisch) und einen 3D-Senkrechtraum (die „hidden layer``) zerlegt; ein
Akzeptanzfenster im Senkrechtraum wählt aus, welche Gitterpunkte in den physischen
Quasikristall projizieren. Krügers exakter Projektor:
\[
P(\tau)=\begin{pmatrix}
1 & -1 & 0 & 0 & \tau & \tau\\
\tau & \tau & 1 & -1 & 0 & 0\\
0 & 0 & \tau & \tau & 1 & -1
\end{pmatrix}\quad(\text{Spalten normiert}),
\]
mit $\tau=\varphi$ (golden) im Parallelraum und $\tau=\sigma=-1/\varphi$ im
Senkrechtraum. Die sechs Spalten sind die sechs ikosaedrischen Fünfachsen (eine
pro Antipodenpaar). Innere Symmetrie: ikosaedrisch (5-zählig); das Punktset ist
aperiodisch. Auffällig für eine Brücke: der versteckte Raum ist \emph{3D} --
dieselbe Dimensionszahl wie FFGFTs Generationenfaser $\mathbb{C}^3$.

% ==============================================================================
\section{Das gemeinsame Objekt: ein Z$_3$-Zirkulant auf der C$_3$-Achse}
% ==============================================================================

Die Ikosaedergruppe enthält C$_3$-Achsen. Die zyklische Koordinaten-Abbildung
$(x,y,z)\mapsto(z,x,y)$ ist eine Drehung der Ordnung~3; sie zerlegt Krügers sechs
Spalten in zwei 3-Zyklen, $\{c_1,c_3,c_5\}$ und $\{c_2,c_4,c_6\}$. Die Gram-Matrix
eines solchen C$_3$-Orbits ist -- durch die C$_3$-Symmetrie erzwungen -- ein
\emph{exakter reeller Z$_3$-Zirkulant}, genau der Typ von FFGFTs Massenoperator
$S=\mathrm{circ}(t_0,t_1,t_2)$ auf $\mathbb{C}^3$:
\[
G_{\text{orbit}}=\mathrm{circ}\!\left(1,\;\tfrac{1}{\sqrt5},\;\tfrac{1}{\sqrt5}\right)
\quad(\tau=\varphi,\ \text{normiert}),
\]
mit $r=2g_1/g_0=2/\sqrt5\approx0{,}894$ und Phase $\theta=0$ (reell). Der
Senkrechtraum ($\tau=\sigma$) gibt $r=-2/\sqrt5$. Geschlossene Form:
$\;|r|=2|\tau|/(1+\tau^2)$, und für $\tau=\varphi$ kürzt sich $\varphi$ exakt zu
$2/\sqrt5$ heraus ($1+\varphi^2=\sqrt5\,\varphi$).

\begin{keyresult}[Kriterien 1 und 2 erfüllt]
Ein Z$_3$-Zirkulant lebt \emph{innerhalb} von HLVs eigener ikosaedrischer Geometrie
-- nicht aufgesetzt, sondern als Gram-Matrix eines C$_3$-Orbits. Das ist dasselbe
Objekt wie FFGFTs Massenoperator (Kriterium~1). Die Abbildung ist explizit
(Kriterium~2): FFGFTs Z$_3$-Generationensymmetrie $\leftrightarrow$ eine C$_3$-Achse
von HLVs Ikosaedergruppe, mit dem Orbit $\{c_1,c_3,c_5\}$ als Bild.
\end{keyresult}

% ==============================================================================
\section{Die Gabelung: 5-zählig gegen 3-zählig}
% ==============================================================================

Der Divergenzpunkt (Kriterium~3) ist scharf und, wie drei Checks zeigen,
\emph{fundamental} -- nicht durch einen Parameter zu schließen.

\textbf{Nur $1/\sqrt5$.} Alle $6\times6$ Skalarprodukte von Krügers Spalten haben
\emph{einen einzigen} Betrag: $1/\sqrt5$. Die 5-zählige Geometrie kennt nur den
Winkel $\arccos(1/\sqrt5)\approx63{,}4^\circ$; den $45^\circ$-Winkel
($\cos=1/\sqrt2$), den FFGFT braucht, gibt es zwischen den Vakuumrichtungen nicht.
Jeder geschützte 3-Orbit liefert also $r=2/\sqrt5$.

\textbf{Kein $\tau$ erreicht $\sqrt2$.} Die geschlossene Form $r(\tau)=2\tau/(1+\tau^2)$
ist für \emph{jedes} $\tau$ durch $1$ gedeckelt (Maximum bei $\tau=1$, kubisch).
FFGFT verlangt $r=\sqrt2\approx1{,}414>1$ -- das liegt \emph{außerhalb} der
erreichbaren Spanne von Krügers Konstruktion, für jeden Parameter.

\textbf{$\sqrt2$ braucht 3-Zähligkeit.} In einer allgemeinen 3-zähligen Geometrie
(zyklische Basis $(a,b,c)$, $r=2(ab+bc+ca)/(a^2+b^2+c^2)$) ist $\sqrt2$ sehr wohl
erreichbar, z.\,B. $v=(1,1,\sqrt2-2^{1/4})$. FFGFTs Koide-Punkt
($\mathrm{circ}(1,\tfrac1{\sqrt2},\tfrac1{\sqrt2})$, Eigenwerte $0{,}29/0{,}29/2{,}41$,
positiv semidefinit) ist also geometrisch legitim -- er verlangt nur eine
\emph{3-zählige} Spalte, nicht Krügers 5-zählige.

\begin{center}
\begin{tabular}{lcccl}
\toprule
 & $r$ & $\theta$ & Koide $Q$ & Quelle \\
\midrule
HLV (geschützt) & $2/\sqrt5=0{,}894$ & $0$ & $7/15$ & 5-zählig / $\varphi$ \\
FFGFT (Massenop.) & $\sqrt2=1{,}414$ & $2/9$ & $2/3$ & 3-zählig / Koide \\
\bottomrule
\end{tabular}
\end{center}

% ==============================================================================
\section{Geschützt gegen generisch: kein Schlupfloch}
% ==============================================================================

Ein Einwand bleibt: der HLV-Quasikristall hat tausende projizierte Punkte --
könnte ein C$_3$-Orbit \emph{irgendeines} Punkts zufällig $\sqrt2$ treffen? Bei
$9822$ projizierten Punkten im Akzeptanzfenster füllen die C$_3$-Orbit-Verhältnisse
$g_1/g_0$ tatsächlich das ganze Kontinuum $[-\tfrac12,1]$, und $1/\sqrt2$ ist
$150$-mal darunter.

\begin{warningboxenv}[Ein generischer Treffer beweist nichts]
Dass generische projizierte Punkte bei \emph{jedem} Winkel sitzen -- auch bei
$1/\sqrt2$ --, gilt für jedes dichte 3-zählige Punktset, sogar für ein kubisches.
Das ist die Stufe-0-Trivialität: man pickt eine \emph{ungeschützte} Richtung mit
dem passenden Winkel. Die \emph{symmetriegeschützten} Orbits -- Krügers sechs
Vakuumrichtungen -- geben dagegen rigide nur $\pm1/\sqrt5$. $\sqrt2$ ist in HLVs
geschützter Struktur nicht vorhanden; es erscheint nur als bedeutungsloser
Zufallswinkel.
\end{warningboxenv}

% ==============================================================================
\section{Die Br\"ucke ist Enthaltensein, kein Widerspruch}
% ==============================================================================
Die obige Gabelung ist eine Aussage \"uber den \emph{gesch\"utzten} Invarianten; sie darf
nicht so gelesen werden, als w\"aren die beiden Modelle unvereinbare Realit\"aten. Auf der
Ebene der Gruppen und Dimensionen sind sie durch \emph{Enthaltensein} vereinbar. Die
Ikosaedergruppe enth\"alt die $3$-Z\"ahligkeit: $C_3<A_5$ -- eine $120^\circ$-Drehung um eine
$(1,1,1)$-Achse und eine $72^\circ$-Vertex-Drehung sind \emph{beide} Symmetrien desselben
Ikosaeders -- und $T^7=T^4\times T^3$ ($7=4+3$). Also $\mathrm{HLV}\supset\mathrm{FFGFT}$:
HLV ist FFGFTs $4$D-Kern, erweitert um die drei Perp-Kreise und die volle $5$-Z\"ahligkeit.
Was die Gabelung also belegt, ist keine Unvertr\"aglichkeit, sondern die \emph{realisierte
gesch\"utzte Geometrie}: FFGFT sch\"utzt eine unabh\"angige $3$-z\"ahlige Spalte ($\sqrt2$,
Koide $2/3$), HLV die $5$-Z\"ahligkeit ($2/\sqrt5$), und $\sqrt2$ ist nicht unter HLVs
gesch\"utzten Richtungen (es erscheint nur als ungesch\"utzter Zufallswinkel, voriger
Abschnitt). Die Modelle teilen das ambiente $3{+}1$ und sind gruppen-/dimensionsm\"a\ss ig
vereinbar, realisieren aber verschiedene gesch\"utzte Struktur -- das ist der empirisch
entscheidbare Unterschied, kein Koordinaten-Artefakt und keine logische Unvertr\"aglichkeit.
Gerechnet in Dok.~285 (\texttt{recompute\_reconcile.py}, \texttt{dim\_bridge\_1..3}).

\medskip\noindent\emph{Wo dies behauptet wird.} Dies alles lebt im kompaktifizierten Bild --- die Tori mit allen Kreisen kompakt, wo kein Kreis als Zeit ausgezeichnet ist. Die Br\"ucke zu physikalischen, dekompaktifizierten Observablen l\"auft \"uber den einen Schritt $S^1\to\mathbb{R}$, den Messvorgang (Dok.~270); das Enthaltensein wird also vor dem Messvorgang, im kompakten Zustand, behauptet --- nicht auf der physikalischen Zeitachse.

% ==============================================================================
\section{HLV ist die dekompaktifizierte Variante: Masse und Zeit nur durch den Zusatz}
% ==============================================================================

Eine Schicht der Gabelung liegt noch tiefer als die Symmetrie-Ordnung und ist
leicht zu übersehen. HLVs Cut-and-Project projiziert in einen unendlichen,
aperiodischen $\mathbb{R}^3$ -- es ist die \emph{dekompaktifizierte} Variante.
FFGFTs Raum ist dagegen \emph{kompakt} ($T^4$). Das ist nicht kosmetisch: Masse
und Zeit sind in FFGFT intrinsisch, weil die Kompaktheit die Moden quantisiert --
das ist die Zeit-Masse-Dualität $\tilde T\cdot m=1$. Ohne Kompaktheit keine
Modenquantisierung, also \emph{kein intrinsischer Massenmaßstab und keine Zeit}.
Genau deshalb ist HLVs geschützte \emph{räumliche} Struktur rein geometrisch: eine
reelle Gram-Metrik mit $\theta=0$ und dimensionslosen Winkeln. Der intrinsische
\emph{Massen}maßstab fehlt, weil die dekompaktifizierte Variante keine
Modenquantisierung hat; er kommt, wie schon in Dok.~282 festgehalten, erst
\emph{durch den Zusatz}. Wichtig ist aber: dieses $\theta=0$ ist die räumliche
Metrik, \emph{nicht} HLVs Zeit. HLV ist nicht zeit- oder phasenlos -- seine Zeit
sitzt in den nativen Spiral-Time-Operatoren (§\,8), die der räumliche Gram gar nicht
erfasst. Der C$_3$-Probe hat nur die Geometrie gesehen, nicht die Dynamik.

Das schärft Kriterium~1. Der gemeinsame Z$_3$-Zirkulant ist nur \emph{algebraisch}
gleich (dieselbe zyklische Form). Physikalisch ist er verschieden: in HLV
(dekompaktifiziert) eine \emph{räumliche Metrik} ohne Masse und Zeit, in FFGFT
(kompakt) der \emph{Massenoperator} mit Masse und Zeit intrinsisch. Die Brücke
trägt die Form, nicht den Inhalt.

\begin{keyresult}[Die Gabelung hat drei Schichten]
\begin{center}
\begin{tabular}{lll}
\toprule
Schicht & HLV & FFGFT \\
\midrule
Kompaktheit & dekompaktifiziert (Metrik) & kompakt ($\tilde T\!\cdot\! m=1$) \\
Symmetrie-Ordnung & 5-zählig ($\varphi$, $1/\sqrt5$) & 3-zählig ($\sqrt2$) \\
Phase & Spiral-Time: schief-adj.\ $R$ (dyn., §\,8) & skalar, $\theta=2/9$ \\
\bottomrule
\end{tabular}
\end{center}
Die oben geforderte komplexe Holonomie ($\theta=2/9$, FFGFT) ist die
\emph{operationale Folge} der Kompaktheit: Masse und Phase entstehen dort durch die
quantisierten Moden. Die HLV-Spalte der Phasen-Zeile ist dabei subtiler als ein
bloßes $\theta=0$ -- das war die \emph{räumliche} Gram-Metrik; HLVs eigene
Zeit-Phase sitzt im schief-adjungierten Spiral-Time-Operator $R$ (§\,8). Eine
vollständige Rekompaktifizierung von HLV muss diese Operatoren \emph{mit-integrieren},
nicht nur den Raum periodisieren -- ein kompaktes HLV ist zudem kein Quasikristall mehr.
\end{keyresult}

% ==============================================================================
\section{Eine kompaktifizierte Variante von HLV: der rationale Approximant}
% ==============================================================================

Lässt sich diese Rekompaktifizierung konkret durchführen, um den Vergleich fair zu
machen? Ja -- mit einem etablierten Begriff der Quasikristall-Physik, dem
\emph{rationalen Approximanten}. Man ersetzt das irrationale $\varphi$ im Projektor
durch rationale Fibonacci-Verhältnisse $\tau_n=F_{n+1}/F_n$
$(1,2,\tfrac32,\tfrac53,\tfrac85,\dots\to\varphi)$. Bei rationalem $\tau=p/q$ wird
der physische Projektor nach Multiplikation mit $q$ \emph{ganzzahlig} -- die
Projektion von $\mathbb{Z}^6$ landet auf einem skalierten Ganzzahlgitter
$(1/q)\mathbb{Z}^3$, also auf einem \emph{periodischen Kristall} auf $T^3$. Bei
irrationalem $\varphi$ gibt es keinen solchen Faktor; das Punktset bleibt dicht und
aperiodisch. Das ist der präzise Sinn von „kompaktifiziertes HLV`` -- und es ist,
wie §\,6 schon sagte, kein Quasikristall mehr, sondern sein periodischer Kristall-Approximant.

Der Approximant ist periodisch, hat also diskrete Moden -- ein \emph{Massenmaßstab
wird möglich}. Damit ist die tiefste Schicht der Gabelung (Kompaktheit)
geschlossen: HLV liegt jetzt im selben Regime wie FFGFT. Das C$_3$-Verhältnis
konvergiert dabei:

\begin{center}
\begin{tabular}{lcc}
\toprule
$\tau=F_{n+1}/F_n$ & $r$ (C$_3$) & Koide $Q$ \\
\midrule
$2/1$ & $0{,}800$ & $0{,}440$ \\
$5/3$ & $0{,}882$ & $0{,}463$ \\
$13/8$ & $0{,}893$ & $0{,}466$ \\
$\to\varphi$ & $0{,}894=2/\sqrt5$ & $7/15$ \\
\bottomrule
\end{tabular}
\end{center}

Aber der Approximant erbt HLVs 5-Zähligkeit: $r\to2/\sqrt5$, Koide $\to7/15$; die
\emph{räumliche} Metrik bleibt reell ($\theta=0$). Damit divergiert die
Symmetrie-Ordnung weiter. Die Phasen-Schicht jedoch ist mit „$\theta=0$`` nicht
erledigt: der rein räumliche Approximant rekompaktifiziert nur den Raum; HLVs
Zeit-Phase liegt in den Spiral-Time-Operatoren, die er noch integrieren muss (§\,8).

\begin{keyresult}[Die Gabelung übersteht den fairen Vergleich]
Mit dem Approximanten ist der Vergleich \emph{kompakt gegen kompakt} und damit
fair. Er bestätigt die Gabelung: \emph{selbst kompaktifiziert} gibt HLV Koide
$7/15$, nicht FFGFTs $2/3$. Die Divergenz war also \emph{kein
Kompakt/Dekompakt-Artefakt} -- sie sitzt echt in der 5- gegen 3-Zähligkeit. Man
kann HLV kompaktifizieren; der Befund hält.
\end{keyresult}

% ==============================================================================
\section{Was die Kompaktifizierung integrieren muss: HLVs Spiral-Time}
% ==============================================================================

Der Approximant aus §\,7 rekompaktifiziert HLVs \emph{Raum}. Aber HLV ist nicht
rein räumlich: seine definierende Dynamik sind die nativen 6D-\emph{Spiral-Time}-Operatoren
\begin{equation}
\Psi(t)=t+i\,\phi(t)+j\,\chi(t),
\end{equation}
mit $\phi(t)$ als Phasen-/Kohärenzkanal und $\chi(t)$ als Gedächtnis-/Hysteresekanal
(Krüger, Zenodo 20668125 und 20512650). Eine kompaktifizierte Variante, die HLV
\emph{bleiben} soll, muss diese Operatoren integrieren -- der rein räumliche
Approximant tut das noch nicht. Das korrigiert die Phasen-Schicht der Gabelung.

Krüger wendet $\Psi$ nicht direkt an (das ergäbe hyperkomplexe Felder), sondern über
eine \emph{konservative Übersetzung} in zwei reelle Operatoren:
\begin{center}
\begin{tabular}{lll}
\toprule
Kanal & wird zu & Anforderung \\
\midrule
$\phi$ (Phase) & Phasen-/Vortizitätsordnung & schief-adjungierter Generator $R$ (energie-neutral) \\
$\chi$ (Memory) & Geschichtsabhängigkeit & positiv-typiger Kern $K$ (dissipativ) \\
\bottomrule
\end{tabular}
\end{center}
mit Exponentialbeispiel $K_\tau(t)=\tfrac1\tau e^{-t/\tau}$, das auf eine interne
Variable reduziert. Daraus folgt dreierlei:

\begin{enumerate}
\item \textbf{Die Phasen-Schicht ist nicht $\theta=0$.} Das $\theta=0$ aus §\,3 ist
die \emph{räumliche} Gram-Metrik, nicht HLVs Zeit. HLVs Phase sitzt im
schief-adjungierten Operator $R$, dynamisch getrieben von $\dot\phi(t)$. Die Schicht
lautet also richtig: \emph{dynamische Operator-Phase $R$} (HLV) gegen \emph{feste
skalare Phase $\theta=2/9$} (FFGFT) -- nicht phasenfrei gegen phasenbehaftet.

\item \textbf{Die Zeit/Gedächtnis-Schicht ist ein echter Brücken-Kandidat.} HLVs
$\chi$-Kanal ist ein positiv-typiger, nicht-Markovscher Gedächtniskern
$\int_0^t K(t-s)\,u(s)\,ds$. FFGFTs Zeit-Erweiterung in Dok.~282 ergibt einen
Gedächtniskern als Fourier-Transformierte der diskreten Spektraldichte
$\sum_k g_k^2\,\delta(\omega-\omega_k)$ -- nicht-Markovsch mit \emph{Rückfluss und
Revivals}, also ein \emph{diskret-spektraler, oszillierender} Kern
(Exponential-Grenzfall $\to$ interne Variable). Das ist nicht ganz dieselbe Gestalt
wie HLVs $\chi$ im Fluid-Sektor, der ein \emph{positiver abklingender} Kern ist
(Krügers Navier-Stokes-Erweiterung, Zenodo~20512650): beide nicht-Markovsch und
positiv-typig, aber verschiedener Form. Wo der Raum forkt (5- gegen 3-zählig), stehen
sich auf der Zeit-Ebene zwei Kerne derselben \emph{Klasse} (nicht-Markovsches
positives Gedächtnis) gegenüber: \emph{räumliche Gabelung, möglicher
$\chi$-Gedächtnis-Brücke}.

\item \textbf{Zwei ehrliche Grenzen.} Erstens: Die $i\!\cdot\!j$-Produktrelation der
Spiral-Algebra ist in \emph{keiner} Quelle Krügers (Dok 271, FA-Hub, Navier-Stokes)
festgelegt -- sie wird stets über die Operator-Übersetzung umgangen. Der
Phasenvergleich ist deshalb \emph{operator-}, nicht algebra-seitig zu führen; ob die
Operator-Phase $R$ FFGFTs $2/9$ trifft, lässt sich erst mit einer expliziten
$R$-Spezifikation entscheiden. Zweitens: Krüger selbst hält fest, dass viele
gewöhnliche reduzierte Modelle Gedächtnis erzeugen; die HLV-Deutung sei nur eine
Kandidaten-Modellklasse, keine demonstrierte Ontologie (Run~M1: HLV-Kern schlägt
schwache Nulls, ist aber auf dem \emph{überdämpften} Beobachtungssektor von einem
generischen Biexponential statistisch \emph{nicht unterscheidbar} -- HLV nur per
Parsimonie/BIC bevorzugt). Die $\chi$-Brücke ist also \emph{strukturelle
Verwandtschaft}, kein Eindeutigkeitsanspruch -- konsistent mit Dok.~276/281 (auch
$\varphi$ ist nicht spezifisch). Das hängt am Sektor: FFGFTs diskret-spektraler Kern
trägt im \emph{Kohärenz-/Revival-Sektor} eine Signatur (BLP-Rückfluss $5{,}125$), die
jeder \emph{abklingende} Kern verfehlt -- die Eindeutigkeit gegen andere
\emph{strukturierte} Kerne bleibt aber offen.
\end{enumerate}

\begin{important}[Die drei Ebenen sind nicht gleich stark]
Liest man die Verbindung nach physikalischer Ebene -- Raum, Zeit, Masse -- so sind
die drei \emph{nicht} gleichwertig, und das ist methodisch entscheidend:
\begin{itemize}
\item \textbf{Masse (Z$_3$-Zirkulant):} die \emph{genuine} Brücke -- aus jeder
Geometrie eigenständig hergeleitet, mit expliziter C$_3$-Abbildung und echtem
Divergenzpunkt. Sie allein besteht das Drei-Teile-Kriterium vollständig.
\item \textbf{Zeit ($\chi$-Gedächtnis):} \emph{strukturelle Verwandtschaft},
Kandidat -- nicht HLV-eindeutig (Run~M1: auf dem überdämpften Sektor ist ein
generischer Biexponential \emph{ununterscheidbar}; HLV nur per Parsimonie bevorzugt).
Bridge-Kandidat, kein Eindeutigkeitsanspruch.
\item \textbf{Raum (Connection-Laplace auf Faserbündel):} \emph{geteilte Sprache},
keine Brücke. Dass beide ein Connection-Laplace auf einem faserwertigen Raum sind,
gilt für \emph{jede} Gitter-Eichtheorie -- das ist Stufe-0-Gerüst, genau das, was das
Kriterium ausschließt. Es darf nicht als tragende Verbindung gezählt werden, sonst
bläht man die Verbindung in Richtung Trivialität auf.
\end{itemize}
Nur die Masse-Ebene ist also eine Brücke im Sinne des Kriteriums; die Zeit-Ebene ein
ehrlicher Kandidat; die Raum-Ebene ein gemeinsamer Rahmen, nicht mehr.
\end{important}

% ==============================================================================
\section{Was übrig bleibt: hexagonal, nicht pentagonal}
% ==============================================================================

Streift man HLV die beiden Achsen ab, an denen es forkt -- Symmetrie-Ordnung und
Kompaktheit --, bleibt genau ein Bestandteil, der nicht schon FFGFT ist: die
Spiral-Time. Ihre Operatoren ($R$, $K$) hängen weder an der 3- noch an der
5-Zähligkeit und überleben den Wechsel auf 3-zählig unverändert. \emph{Was von HLV
substanziell bleibt, ist also nicht eine Geometrie, sondern eine Zeit} -- und genau
dort, in der Dynamik, trifft sie FFGFTs nicht-Markovschen Kern aus Dok.~282. Der
geometrische Rest -- die ikosaedrische 5-Zähligkeit -- ist ein eleganter, aber
physikalisch leerer Ast: kein Sektor realisiert ihn, und $\varphi$ ist nicht
spezifisch (Dok.~276/281).

Daraus folgt eine knappe geometrische Lehre. Die 5-Zähligkeit ist
\emph{nicht-kristallographisch} -- sie parkettiert nicht periodisch und erzwingt die
Dekompaktifizierung; die 3-Zähligkeit ist \emph{kristallographisch}, aus derselben
Familie wie die 6-Zähligkeit (Dreiecks-/Hexagonalgitter). Wo Geometrie die
Massenstruktur trägt, ist sie damit hexagonal, nicht pentagonal -- und das ist kein
neuer Befund, sondern FFGFTs eigener Anker: das \emph{Euler-Tonnetz} ist ein
hexagonales Gitter (Dok.~275, „der Rahmen ankert im Euler-Tonnetz, nicht im Goldenen
Schnitt``).

\begin{important}[Die belastbare Form der These]
„Hexagonal`` meint die kristallographische Familie (3-/6-zählig), nicht die
Behauptung, die Raumzeit sei eine Hexagonal-Parkettierung -- FFGFTs Basis bleibt
$T^4$ mit $Z_3$-Generationensymmetrie. Empirisch 3-zählig bestätigt sind bisher
allein die geladenen Leptonen; die Hadronen sind fast entartet ($Q\approx1/3$), also
weder 3- noch 5-zählig. Belastbar bleibt: \emph{wo Geometrie die Massenstruktur
trägt, ist sie hexagonal/kristallographisch; das Pentagonale bleibt ein eleganter,
physikalisch leerer Ast; und was von HLV substanziell übrig bleibt, ist seine Zeit,
nicht seine Geometrie.}
\end{important}

% ==============================================================================
\section{Fazit}
% ==============================================================================

Alle drei Teile des Kriteriums sind beantwortet. (1)~Gleiches Objekt: ein
Z$_3$-Zirkulant aus HLVs eigener Geometrie. (2)~Abbildung: die C$_3$-Achseneinbettung
von FFGFTs Z$_3$ in HLVs Ikosaedergruppe. (3)~Divergenz: fundamental, und in drei
Schichten -- Kompaktheit (dekompaktifiziert/Metrik gegen kompakt/$\tilde T\!\cdot\! m=1$),
Symmetrie-Ordnung (5-zählig $\varphi$, $1/\sqrt5$ gegen 3-zählig Koide, $\sqrt2$)
und Phase (HLVs dynamische Spiral-Time-Operator-Phase $R$ gegen FFGFTs feste
skalare $2/9$; §\,8). Die Kompaktheit bleibt entscheidend für den \emph{Massen}maßstab
(ohne quantisierte Moden kein intrinsisches $m$); HLVs \emph{Zeit} jedoch sitzt in
der Spiral-Time, unabhängig davon. Kein Parameter von HLV erreicht FFGFTs geschützten
Koide-Punkt -- doch auf der Zeit/Gedächtnis-Ebene stehen sich beide nicht-Markovschen
Kerne derselben Klasse gegenüber: \emph{räumliche Gabelung, möglicher Zeit-Brücke} (§\,8).

Das ist genau, was das Kriterium liefern soll: eine \emph{genuine} Verbindung --
gemeinsame Z$_3$-Zirkulant-Familie, explizite C$_3$-Abbildung -- \emph{mit} einem
echten Divergenzpunkt, und nachweislich \emph{keine} Stufe-0-Trivialität (das
Geschützte trennt sich, nur das Bedeutungslose fällt zusammen). FFGFT und HLV
sitzen an zwei \emph{verschiedenen} Punkten derselben Familie. Es ist eine
\emph{Gabelung}, kein Zusammenfall: gemeinsamer Stamm (Z$_3$-Zirkulant aus
höherdimensionalem Gitter mit verstecktem Raum), zwei Äste je nach Symmetrie-Ordnung.

Für die Korrespondenz ist die Aussage präzise und fair: \emph{Wir treffen uns im
Z$_3$-Zirkulanten entlang deiner C$_3$-Achse; wir trennen uns, weil deine
ikosaedrische 5-Zähligkeit rigide $1/\sqrt5$ erzwingt, während FFGFTs
Koide-Bedingung 3-zählig $\sqrt2$ verlangt -- und kein Parameter deiner
Konstruktion überbrückt das.} Reproduzierbar mit
\texttt{ffgft\_hlv\_c3\_bridge\_probe.py} (sechs Checks, numpy-only).
